A linear probe that clears significance is routinely read as evidence
that a model has learned, represents, or is developing a capability.
The inference is unsound: a network's activations are a
high-dimensional expansion of the visible tokens, and a linear readout
on such an expansion decodes many functions of the input that training
never put there. This paper prescribes three controls that make the
inference testable, an untrained-weights twin of the model run through
the identical probe pipeline as an acceptance test, basis-starved
splits whose validation items share no surface-component values with
probe training, and tolerances calibrated from the mechanism that
generates each control's events, and reports two preregistered probing
campaigns that the controls terminated before any outcome measurement.
In the first, the untrained twin fired on 120 of 120 fits: under
standard splits, probe significance measured reservoir decodability,
not learning. In the second, starving closed the lookup class and the
twin exposed a class beneath it, 13 of 25 capabilities decodable from
untrained weights at margins of .06 to .82 via surface statistics that
no value holdout can remove; the twelve survivors read exactly zero
untrained margin in every cell. A successor battery built with the
screen at inclusion time froze with the untrained gate clean, the
program's first, and then committed the outcome-side form of the same
error: an untrained floor sits at zero on a generation task, so it
credited format acquisition as capability, and re-scoring against
within-answer-space chance moved the headline rank correlation from
.368 to .200. The two capabilities carrying the second and third
highest starved margins scored exactly zero argmax accuracy at every
scale: decodable is not learned, and on this evidence not generable
either. A taxonomy of the leak mechanisms and a checklist written for
verbatim adoption distill the record.
